\chapter{Unitary representations}
\setcounter{exercisepar}{0}

For every vector space E, we denote by \(1_{E}\) the identity map of E. If E, F and G are vector spaces and if u: \(F \rightarrow E\) and v: \(G \rightarrow F\) are linear maps, one sometimes denotes by uv the linear map \(u \circ v\) from G to E.

For a Hilbert space E, one denotes by \(\langle x | y \rangle_{E}\) , or simply \(\langle x | y \rangle\) , the scalar product on E.

If X is a locally compact topological space and \(\mu\) a measure on X, we denote by \(\mathcal{L}^{p}(\mathrm{X},\mu)=\mathcal{L}_{\mathbf{C}}^{p}(\mathrm{X},\mu)\) and \(\mathrm{L}^{p}(\mathrm{X},\mu)=\mathrm{L}_{\mathbf{C}}^{p}(\mathrm{X},\mu)\) for \(p\in[1,+\infty]\) .

One denotes by e the identity element of a group G.

Let A be an associative algebra, E a topological vector space and \(\pi\colon\mathrm{A}\to\mathcal{L}(\mathrm{E})\) a representation of A in E. One says that \(\pi\) is non-degenerate if the set of vectors \(\pi(a)x\) for \(a\in\mathrm{A}\) and \(x\in\mathrm{E}\) is total in E.

